\begin{afterword}
\summaryhead

The post tells the story of group selection. Before 1966, biologists such as Wynne-Edwards held that animals restrain their breeding so as not to exhaust their food, and explained this by selection among isolated subpopulations: groups that restrained themselves would survive and recolonize. The author grants that this is possible but says it needs a low cost of restraint, a high benefit and very small groups, conditions that a simulation suggests are rare. The group selectionists lost the argument, and observation went against them.

Then Michael Wade selected laboratory populations of flour beetles for small size. The beetles did not restrain their breeding; they evolved to eat eggs and larvae. The author concludes that the group selectionists started from an aesthetic wish for nature in harmony, found a mechanism to justify it, and never asked neutrally what that mechanism would produce. Evolutionary biology, the post ends, is good training against this kind of rationalization.

\responsehead

The biology in the post is sound. It states the group-selection mechanism fairly, grants that it is possible, and explains the free-rider problem that makes it hard. It describes the design of Wade's experiment accurately. The problems are in the history and the psychology.

The post's lesson depends on a claim about how the group selectionists reasoned: that they began from ``the beautiful idea'' of harmony, searched for a reason, and stopped. The post gives no evidence for this. It quotes none of them and cites no history. The one first-hand account I found, Wynne-Edwards's own note of 1980, says the idea came to him while reviewing a book on animal numbers; it neither confirms nor rules out the post's story. The post's evidence is the laboratory result, which bears on what group selection produces, not on how anyone arrived at a prediction. The only evidence it gives that the result was foreseeable is that it is ``massively obvious in retrospect,'' the judgement the author's own ``Hindsight Devalues Science'' warns against. Better support existed: cannibalism was a known regulator of flour beetle numbers before the experiment.

The post also makes its opponents' view gentler than it was. Wynne-Edwards did not describe voluntary restraint. He described territories and peck-orders that exclude animals from feeding and breeding, and surplus animals that emigrate or die; in red grouse, he wrote, 60 per cent of the birds become outcasts and die in an average year.

Two facts about Wade's work are missing, and a reader should know them. Cannibalism was one of four traits that changed; the abstract lists fecundity first. And the experiment is also cited for showing group selection to be strong: the response was large and fast, and Wade concluded with Slatkin that group selection ``can be much more effective'' in nature than is commonly supposed. The wider question is still disputed, as the 2010 exchange in \textsc{Nature} between Nowak, Tarnita and Wilson and 137 critics shows; the post's own statement of the revival is compatible with both sides.

In short: an accurate account of why group selection is hard, combined with an unsourced story about how the group selectionists reasoned, supported mainly by hindsight, the error the author warned against in ``Hindsight Devalues Science.''
\end{afterword}
